\section{Related Work}
\label{sec:related}

\paragraph{Selecting items from historical results.}
Anchor Points~\citep{vivek2024anchor} clusters items by how a pool of source models responds to them, keeps the cluster medoids as anchors and estimates a target's score as a cluster-size-weighted average over the anchors.
tinyBenchmarks~\citep{polo2024tiny} combines such anchors with item-response-theory (IRT) models fitted on the source models, and releases reduced benchmarks;
metabench~\citep{kipnis2025metabench} selects items by IRT information and predicts with a generalised additive model;
\citet{pacchiardi2024100} train an ``assessor'' on the target's results over a reference set, and \citet{li2025active} learn an adaptive item-acquisition policy; \citet{balkir2026confident} adapt IRT-based adaptive testing to continuous scores and ranking.
These works use different model pools, splits (random, by date, by organisation), metrics and numbers of repetitions, so their reported errors are not comparable with each other;
our contribution on this axis is a single protocol over the same matrices, with the same budget, source/target split and metrics for every method.

\paragraph{The closest comparison.}
\citet{zhang2025benchmark} compare eleven benchmark-prediction methods over many benchmarks and release the score matrices and implementations that we use (Section~\ref{sec:protocol}).
Their repository states conclusions about which methods are competitive and which fail when targets are stronger than the source models; we treat those statements as hypotheses to be tested, not as inputs.
We re-implement all methods ourselves (the released \texttt{benchpred} package is inspected but not used), remove exact duplicate model rows before splitting, add a family-held-out split, report model-level bootstrap intervals and paired differences, and add a model-assisted estimator and its ablations.

\paragraph{Model-assisted estimation and prediction-powered inference.}
Estimating a finite-population mean from a probability sample, using a model fitted elsewhere only to reduce variance while the design supplies validity, is the model-assisted (difference-estimator) approach of survey sampling~\citep{sarndal1992model,breidt2017model}.
Prediction-powered inference~\citep{angelopoulos2023ppi} is a closely related construction for settings in which model predictions are available on many items and labels on few, and stratified variants have been used to evaluate language models~\citep{fisch2024stratified}.
In our setting the ``population'' is the benchmark's item set, the ``labels'' are the target model's 0/1 outcomes on the sampled items, and the ``predictions'' come from a latent-trait model fitted on the source models.
We do not claim the estimator family as new; what we study is how its design (stratification, probability-proportional-to-size allocation) and predictor behave under the protocols above, in particular when targets lie above the range of the source models.
Related uses of sparse or subsampled evaluation include \citet{ghosh2024onebench}, who aggregate sparse measurements across datasets, and \citet{polo2024prompteval}, who estimate performance distributions over prompt variants.
